\scanpage{114}
de cardinal $c$, donc pour des raisons de cardinal on aura

\[ \widetilde E_c(r_1)=\widetilde E_c(r'_2) \]

Supposons $\underline{c\geq3}$, de sorte que pour $r_1,r_2\in A'$ donnés, \inserted{$r_1\neq r_2$} $\exists\,r_3\in A^{\unsure{\prime|\prime\prime}}$ tel que $r_3\neq r_1,\deleted{\unsure{r_2'}}\unsure{r_2|r'_2}$. On trouve donc $\widetilde E_c(r_1)=\widetilde E_c(r_2)\inserted{=\widetilde E_c(r'_3)}$, donc tous les $\widetilde E_c(r_i)$ \deleted{$(r_i\in A')$} \deleted{sont égaux, et $\widetilde E_c(r_i')$ \unsure{eux égaux|sont égaux}} sont égaux et égaux aux $\widetilde E_c(r'_i)$, or $\widetilde E_c(r_i')=\complement\,\widetilde E_c(r_i)$, ce qui est absurde.

\uline{Lemme 2} Soient \deleted{$\unsure{v_1,v_2}\inserted{\overline{v_1,v_2}}\in\widetilde E_b,\ u\in\widetilde E_b$ lié à $s$ et $t$,} $\underline t_1=(u,v_1,w_1)$, $\inserted{\underline t_2=}(u,v_2,w_2)$ deux triangles sup. sur $(a,b,c)$, ayant le sommet $u$ (sur $a$) en commun. Considérons $\pi=\sigma_{\unsure{v_1|w_1}}\sigma_{\unsure{v_2|w_2}}$ (autom \inserted{involutif} qui fixe $a,b,c$) alors $\pi$ \deleted{échange les deux triangles précédents, i.e. $\pi u=u,\ \pi v_1=w_2,\ \pi w_1=w_2$} \inserted{\deleted{\unsure{ne}} échange $w_1$ et $w_2$}

\uline{Corollaire} \inserted{(Si $c\neq1,2$)} \deleted{Soient} $a,b$ deux sommets distincts de $\underline t_0$, $G_b$ le \inserted{ss} groupe \inserted{du grpe} des autom de $(S,A)$ engendré par les produits de \deleted{deux} \inserted{$2$} \deleted{symétries} \unsure{forme|par} $\sigma_{v_1}\sigma_{v_2}$, avec $v_1,v_2\in E_b$ \unsure{ou} (donc $G_b$ est un \deleted{\unsure{2-grpe com}} $\mathbb F_2$-vectoriel$\unsure{\subseteq|G}$, et il \unsure{invarie|inverse|respecte} $a,b,c$) Alors $G_b$ est transitif sur $\widetilde E_a$, de façon plus précise, si

% archive sheet number: 113
